\section{Introduction}

The goal of implicit computational complexity (ICC) is to provide a programming language which characterizes a given complexity class. Concretely, this means that a procedure is representable in the language if and only if it inhabits the desired complexity class. ICC has been extensively studied for many complexity classes including P \cite{hofmannLinearTypesNonsizeincreasing1999,hofmannTypeSystemBounded2000,hofmannLinearTypesNonsizeincreasing2003,gaboardiSoftLinearLogic2008}, PSPACE \cite{hofmannLinearTypesNonsizeincreasing2003, gaboardiSoftLinearLogic2008,gaboardiImplicitCharacterizationPSPACE2012}, EXPTIME \cite{hofmannStrengthNonsizeIncreasing2002}, and LOGSPACE \cite{schoppStratifiedBoundedAffine2007, ramyaaRamifiedCorecurrenceLogspace2011}. Such languages often achieve their desired complexity class via a slight modification to a standard linear or affine type system \cite{dallagoImplicitComputationComplexity2022}.

One prominent such programming language for the complexity classes P and FP is Martin Hofmann's Linear Function Programming Language (LFPL), first presented in \cite{hofmannLinearTypesNonsizeincreasing1999}, where it is shown that LFPL is polynomial-time sound. That is, all functions definable in LFPL belong to FP. The logical inverse of this statement, polynomial-time completeness, is later proven in \cite{hofmannStrengthNonsizeIncreasing2002}. LFPL is notable for characterizing FP while also supporting a very natural programming style. Recently, there has been some renewed interest in LFPL. Atkey \cite{atkeyPolynomialTimeDependent2024} extends LFPL with dependent types and unrestricted computation at the type level, creating a language equipped with the powerful guaranteees of dependent type theory while still enforcing that computation is polynomial-time at the term level. Lorenzen et al. \cite{lorenzenFP2FullyInPlace2023} utilizes the type system of LFPL in their language Koka to ensure that certain functional programs can be executed fully in-place.

Despite its prominence, we are not aware of a fully self-contained presentation of LFPL and its core metatheory. The original presentation by Hofmann \cite{hofmannLinearTypesNonsizeincreasing1999} contains a semantic soundness proof but and only a weak completeness result for linear-space polynomial-time Turing machines. A full but somewhat informal completeness result is given as a corollary in \cite{hofmannStrengthNonsizeIncreasing2002}, but polynomial-time is not the main focus. Aehlig et al. \cite{aehligSyntacticalAnalysisNonsizeincreasing2002} provides a syntactic soundness proof for a small-step semantics and a weaker completeness proof for a strong extension of LFPL. Atkey \cite{atkeyPolynomialTimeDependent2024} presents a similar completeness proof, as well as a soundness proof for LFPL extended with dependent type theory. Lorenzen et al. \cite {lorenzenFP2FullyInPlace2023} presents LFPL's type system but no metatheory. None of these works contain both a soundness and completeness proof of the original LFPL system, and moreover we are not aware of any fully correct completeness proofs.

In this article, we give a modern account and mechanization of LFPL and its metatheory. Our primary goal is for this article to be a self-contained and accessible resource for those who are interested in learning about LFPL, with all of our results guaranteed correct by the mechanization. We have witnessed firsthand the necessity of such a resource. It was challenging to collect all this information on LFPL from across the various different papers we have mentioned, and during the process we discovered some errors in the completeness proof presented in \cite{hofmannStrengthNonsizeIncreasing2002}, which is the only full completeness proof of the original LFPL that we are aware of.

Our presentation of LFPL involves a denotational semantics similar to the original given in \cite{hofmannLinearTypesNonsizeincreasing1999} as well as a big-step operational cost semantics, which to our knowledge had not yet been developed for LFPL. One advantage of these two semantics is that they both avoid making use of syntactic substitution, which is a notoriously difficult operation to support and reason about in proof assistants. By using such a semantics for LFPL, we greatly reduce the complexity of our mechanization.

The mechanized completeness proof is a simplified and streamlined version of the one given by Hofmann \cite{hofmannStrengthNonsizeIncreasing2002}. The main novelty is the implementation of a stack-like data structure in LFPL that allows us to avoid much of the complexity of Hofmann's proof as well as sidestep the errors we discovered in that proof. To our knowledge, ours is the strongest fully correct completeness proof for LFPL. Hofmann gives a weaker completeness theorem for linear-time Turing machines only in \cite{hofmannLinearTypesNonsizeincreasing1999}, and \cite{aehligSyntacticalAnalysisNonsizeincreasing2002, atkeyPolynomialTimeDependent2024} both contain full completeness proofs but for powerful extensions of LFPL, thereby weakening the completeness result. Our completeness proof works with a minimal version of LFPL much like the one presented in \cite{hofmannLinearTypesNonsizeincreasing1999}, using the fewest tools to simulate a polynomial-time Turing machine.

The mechanized soundness proof is inspired by Aehlig et al \cite{aehligSyntacticalAnalysisNonsizeincreasing2002}. We give a modern presentation of their result for a big-step operational cost semantics. Using such a semantics allows us to reason about concrete program executions and avoids the substitution-related challenges that arise in mechanizing a small-step semantics. While a big-step operational cost semantics is farther away from a concrete machine model, it is well-known how to link it to a small-step cost semantics, which have been shown to be reasonable with respect to lower-level semantic models (see \cite{blellochProvableTimeSpace1996}). To our knowledge, ours is the strongest soundness proof for LFPL; we strengthen it by extending the language with several powerful features such as a built-in stack data structure and an exponential modality.

In summary, we make the following contributions: The first complete presentation of the LFPL metatheory, a novel completeness proof, a novel soundness proof, and the first complete mechanization of the LFPL metatheory. These novelties make our proofs accessible and amenable to mechanization; our desire to mechanize the LFPL metatheory is partially responsible for driving us to discover these proofs. Our completeness proof resolves some errors we discovered in the only other completeness proof of equal strength to ours.